\chapter{الجبر الخطي}\label{ch:b2:linalg}

كان \omterm{def:b2:structures:algebra}{الجبر} الخطي في مجلد السنة الأولى يعمل على $\R$ أو $\C$ في
البعد المنتهي، ويقبل \omterm{def:b2:linalg:det}{المحدد} العام دون برهان. ويرفع هذا الفصل
القيود الثلاثة جميعًا: فتُصاغ النظرية على حقل كيفي $K$،
ويُطوَّر بانتظام التفاعل بين فضاء و\emph{ثنويّه}
(الأسس الثنوية، والمُبيدات، والمنقولات)، و\emph{يُبنى} \omterm{def:b2:linalg:det}{المحدد} أخيرًا
انطلاقًا من الصيغ المتعددة الخطية المتناوبة ومن إشارة
\cref{ch:b2:structures} --- فيُسدَّد بذلك كل ما قُبل دون برهان في السنة الأولى.

وفي كل ما يلي، $K$ حقل ($\Q$ أو $\R$ أو $\C$ أو $\Z/p\Z$ --- فالنظرية
لا تعبأ بذلك)، والفضاءات، ما لم يُذكر خلاف ذلك، منتهية البعد
على $K$. ونتائج السنة الأولى (الأسس، والبعد،
ومبرهنة الرتبة، والمصفوفات) تنتقل حرفيًا: فبراهينها لم تستعمل قط
شيئًا غير بديهيات الحقل.

\section{الفضاء الثنوي}

\begin{definition}[الفضاء الثنوي، الأساس الثنوي]\label{def:b2:linalg:dual}
\emph{الثنوي}\index{فضاء ثنوي} للفضاء $E$ هو $E^* = \mathcal{L}(E,
K)$، أي فضاء الصيغ الخطية. وإذا كانت $\mathcal{B} = (e_1, \dots, e_n)$
أساسًا للفضاء $E$، فإن \emph{صيغ الإحداثيات} $e_1^*, \dots, e_n^*$
المعرَّفة بالعلاقة $e_i^*(e_j) = \delta_{ij}$ (رمز كرونيكر: $1$ إذا كان $i = j$،
و$0$ فيما عدا ذلك) تشكّل \emph{الأساس الثنوي}\index{أساس ثنوي}
$\mathcal{B}^*$ للفضاء $E^*$؛ وعلى الخصوص $\dim E^* = \dim E$، و
\[
x = \sum_{i=1}^{n} e_i^*(x)\, e_i \quad (x \in E),
\qquad
\varphi = \sum_{i=1}^{n} \varphi(e_i)\, e_i^* \quad (\varphi \in
E^*).
\]
\end{definition}

\begin{proof}[برهان أن $\mathcal{B}^*$ أساس]
حرّ: بتطبيق تركيب خطي معدوم $\sum \lambda_i e_i^* = 0$ على
$e_j$ نجد $\lambda_j = 0$. مولّد: من أجل $\varphi \in E^*$، تُعدم
الصيغة $\varphi - \sum_i \varphi(e_i) e_i^*$ كل $e_j$، ومن ثم فهي معدومة
(فالتطبيق الخطي المنعدم على أساس منعدم). والصيغتان
المعروضتان هما الحسابان أنفسهما مقروءين في الاتجاه المباشر.
\end{proof}

\begin{example}\label{ex:b2:linalg:dualexamples}
على $K_n[X]$ بالأساس $(1, X, \dots, X^n)$: \omterm{def:b2:linalg:dual}{الأساس الثنوي} هو $P
\mapsto \frac{P^{(k)}(0)}{k!}$ (أي معاملات تايلور). وثمة أساس آخر
للثنوي: التقييمات $P \mapsto P(x_i)$ عند $n + 1$ نقطة مختلفة --- وأساسه
``القبلي'' في $K_n[X]$ هو بالضبط عائلة
كثيرات حدود لاغرانج $L_i$ (مجلد السنة الأولى)، لأن $L_i(x_j) =
\delta_{ij}$. فالاستكمال \emph{هو} الثنوية.
\end{example}

\begin{method}[الأساس الثنوي والأساس القَبْلي عمليًا]\label{met:b2:linalg:dualbasis}
لنشر صيغة $\varphi$ على أساس $(e_i)$ للفضاء $E$:
الإحداثيات هي \emph{القيم} $\varphi(e_i)$ --- ولا جملة تُحلّ.
ولإيجاد الأساس $(u_j)$ للفضاء $E$ الذي ثنويّه أساس معطى
$(\varphi_1, \dots, \varphi_n)$ للفضاء $E^*$ (وهو
\emph{الأساس القبلي}): نحلّ الجمل الخطية $n$
\[
\varphi_i(u_j) = \delta_{ij} \qquad (1 \leq i \leq n),
\]
عمودًا $u_j$ في كل مرة؛ وبلغة المصفوفات، إذا كانت أسطر $M$
تسرد معاملات العناصر $\varphi_i$ في أساس معلوم للفضاء
$E^*$، فإن أعمدة $M^{-1}$ هي $u_j$. أما وجود الأساس القبلي
ووحدانيته فيُبرهن عليهما في مسألة نهاية الأسبوع في هذا الفصل؛
والحساب هو هذا القلب دائمًا.
\end{method}

\begin{example}[أساس ثنوي للفضاء $\R^2$، محسوب بالكامل]\label{ex:b2:linalg:dualbasis2d}
من أجل الأساس $b_1 = (1, 1)$، $b_2 = (1, -1)$ للفضاء $\R^2$: يجب أن يحقق
\omterm{def:b2:linalg:dual}{الأساس الثنوي} $(b_1^*, b_2^*)$ الشرط $b_i^*(b_j) =
\delta_{ij}$. وبكتابة $b_1^*(x, y) = \alpha x + \beta y$، يعطي
الشرطان $\alpha + \beta = 1$ و$\alpha - \beta = 0$
\[
b_1^*(x, y) = \frac{x + y}{2},
\qquad\text{وكذلك}\qquad
b_2^*(x, y) = \frac{x - y}{2} .
\]
وللتحقق: $b_1^*$ \emph{ليست} $e_1^* + e_2^*$ مقيَّمة
ببساطة --- \omterm{def:b2:linalg:dual}{فالأساس الثنوي} يتعلق بالأساس كله، لا بكل
متجهة على حدة (فتعويض $b_2$ بالمتجهة $(0, 1)$ يحوّل
$b_1^*$ إلى $x \mapsto x$). وصيغة النشر تعمل:
$(x, y) = \frac{x+y}2\,b_1 + \frac{x-y}2\,b_2$، وهو تفكيك
الزوج إلى جزء زوجي وآخر فردي --- فالأسس الثنوية مستخرِجات
إحداثيات، وهذا الأساس يستخرج الجزأين المتماثل وضدّ المتماثل.
\end{example}

\begin{definition}[المُبيد]\label{def:b2:linalg:annihilator}
من أجل فضاء جزئي $F \subseteq E$، يكون
\emph{المُبيد}\index{مبيد} هو
\[
F^{\circ} = \{\varphi \in E^* : \varphi|_F = 0\},
\]
وهو فضاء جزئي من $E^*$.
\end{definition}

\begin{theorem}[بعد المُبيد]\label{thm:b2:linalg:annihilator}
$\dim F^{\circ} = \dim E - \dim F$. زيادة على ذلك، يعكس التطبيق $F \mapsto F^\circ$
الاحتواءات، ويُستعاد $F$ من مُبيده:
\[
F = \{x \in E : \forall\varphi \in F^\circ,\ \varphi(x) = 0\}.
\]
ومن ثم فكل فضاء جزئي بعده $p$ في بعد $n$ هو
مجموعة حلول $n - p$ معادلة خطية مستقلة --- والعكس بالعكس.
\end{theorem}

\begin{proof}
لنختر أساسًا $(e_1, \dots, e_p)$ للفضاء $F$ ونُتمّه إلى أساس للفضاء
$E$. تُبيد الصيغة $\varphi = \sum \varphi(e_i) e_i^*$ الفضاء $F$ إذا وفقط إذا انعدمت
معاملاتها $p$ الأولى: أي $F^\circ =
\operatorname{Vect}(e_{p+1}^*, \dots, e_n^*)$، وبعده $n - p$.
وعكس الاحتواءات مباشر. وأما الاستعادة: فالطرف الأيمن
يحتوي $F$؛ وعكسيًا، إذا كان $x \notin F$، فنُتمّ أساسًا للفضاء
$F$ بالمتجهة $x$ ثم بمتجهات أخرى؛ فصيغة إحداثيات $x$ في هذا
الأساس تُبيد $F$ ولا تُبيد $x$. وأما القراءة بلغة ``المعادلات'' فتأخذ
أساسًا $(\varphi_1, \dots, \varphi_{n-p})$ للفضاء $F^\circ$: فيكون عندئذٍ $F =
\bigcap \ker\varphi_j$، وهو تقاطع $n - p$ فضاء
فوقي مستقلًا.
\end{proof}

\begin{example}[مُبيد، في الاتجاهين]\label{ex:b2:linalg:annihilatorwork}
ليكن $F = \operatorname{Vect}\bigl((1, 2, 1),\ (1, 0, -1)\bigr)
\subseteq \R^3$. تُبيد الصيغة $\varphi = a\,e_1^* + b\,e_2^* + c\,e_3^*$
الفضاء $F$ إذا وفقط إذا
\[
a + 2b + c = 0
\qquad\text{و}\qquad
a - c = 0 ,
\]
أي $c = a$ و$b = -a$: ومنه $F^\circ = \R\,(e_1^* - e_2^* +
e_3^*)$، وبعده $3 - 2 = 1$ كما
تقتضي \cref{thm:b2:linalg:annihilator}. وبقراءتها في الاتجاه المعاكس:
$F = \{(x, y, z) : x - y + z = 0\}$ --- فقد استُعيد المستوي بوصفه
نواة الصيغة الوحيدة التي تولّد $F^\circ$. فالانتقال من عائلة
مولّدة إلى معادلات \emph{هو} حساب مُبيد؛
والانتقال من المعادلات إلى تمثيل وسيطي هو حساب
مُبيد قَبْلي. (وللتحقق: المتجهتان المولّدتان تحققان $x - y +
z = 0$.)
\end{example}

\begin{definition}[التطبيق المنقول]\label{def:b2:linalg:transpose}
من أجل $u \in \mathcal{L}(E, F)$، يكون \emph{المنقول}\index{تطبيق
منقول} $u^{\mathsf T} \in \mathcal{L}(F^*, E^*)$ هو
\[
u^{\mathsf T}(\psi) = \psi \circ u .
\]
وهو يحقق $(v \circ u)^{\mathsf T} = u^{\mathsf T} \circ v^{\mathsf
T}$، ومصفوفة $u^{\mathsf T}$ في الأسس الثنوية هي
منقول مصفوفة $u$ --- وهذا ما \emph{يفسّر} أخيرًا
منقول السنة الأولى.
\end{definition}

\begin{example}[المنقول، عنصرًا عنصرًا]\label{ex:b2:linalg:transposework}
ليكن للتطبيق $u \colon \R^2 \to \R^3$ المصفوفة $A =
\begin{psmallmatrix} 1 & 2\\ 0 & 1\\ 3 & 0\end{psmallmatrix}$ في
الأسس القانونية. من أجل $\psi = b_1f_1^* + b_2f_2^* + b_3f_3^*
\in (\R^3)^*$، نحسب $u^{\mathsf T}(\psi) = \psi \circ u$ على
أساس الفضاء $\R^2$:
\[
(\psi \circ u)(e_1) = \psi(1, 0, 3) = b_1 + 3b_3,
\qquad
(\psi \circ u)(e_2) = \psi(2, 1, 0) = 2b_1 + b_2 .
\]
ومنه $u^{\mathsf T}(\psi) = (b_1 + 3b_3)\,e_1^* + (2b_1 +
b_2)\,e_2^*$، وتكون مصفوفة $u^{\mathsf T}$ في الأسس الثنوية هي
\[
\begin{pmatrix} 1 & 0 & 3\\ 2 & 1 & 0\end{pmatrix}
= A^{\mathsf T} :
\]
\omterm{def:b2:linalg:transpose}{فالمنقول} المجرد \emph{هو} المصفوفة المقلوبة، ولم يبقَ حساب
يُؤخذ على الثقة. ولاحظ الآلية: صار \emph{العمود} رقم
$j$ من $A$ هو \emph{السطر} رقم $j$ في المصفوفة
الجديدة لأن $\psi \circ u$ يقرأ مخرجات $u$ عبر
معاملات $\psi$.
\end{example}

\begin{proposition}\label{prop:b2:linalg:transposerank}
$\ker u^{\mathsf T} = (\operatorname{im} u)^{\circ}$ و
$\operatorname{im} u^{\mathsf T} = (\ker u)^{\circ}$. ومن ثم
$\operatorname{rk}(u^{\mathsf T}) = \operatorname{rk}(u)$: أي إن رتبة الأسطر
تساوي رتبة الأعمدة، وقد بُرهن على ذلك بنيويًا.
\end{proposition}

\begin{proof}
تُعدم $\psi \in \ker u^{\mathsf T} \iff \psi \circ u = 0 \iff \psi$ الفضاء
$\operatorname{im} u$: وهذه هي المتطابقة الأولى. وأما الثانية: فتُعدم $u^{\mathsf
T}(\psi) = \psi \circ u$ الفضاء $\ker u$ دائمًا، ومنه
$\operatorname{im} u^{\mathsf T} \subseteq (\ker u)^\circ$؛
وتتطابق الأبعاد بمبرهنة الرتبة وبالعلاقة
\cref{thm:b2:linalg:annihilator}:
\[
\operatorname{rk} u^{\mathsf T} = \dim F^* - \dim\ker u^{\mathsf T}
= \dim F - \bigl(\dim F - \operatorname{rk} u\bigr)
= \operatorname{rk} u
= \dim (\ker u)^{\circ} . \qedhere
\]
\end{proof}

\begin{example}[الرتبة مقروءة من الجهتين]\label{ex:b2:linalg:rankbothsides}
ليكن
\[
A = \begin{pmatrix}
1 & 2 & 0 & 1\\
0 & 1 & 1 & 1\\
1 & 3 & 1 & 2
\end{pmatrix} .
\]
\emph{رتبة الأعمدة:} السطر الثالث مجموع الأولين، ومنه
$\operatorname{rk} A \leq 2$؛ والعمودان $1$ و$2$ حرّان:
ومنه $\operatorname{rk} A = 2$. \emph{نواة \omterm{def:b2:linalg:transpose}{المنقول}:} يعطي حلّ
$A^{\mathsf T}y = 0$ أن $y \in \R\,(1, 1, -1)$، ومنه فبعد $\ker
A^{\mathsf T}$ هو $1 = 3 - 2$: أي بالضبط
$(\operatorname{im} A)^\circ$ بمماهاة
$(\R^3)^*$ مع متجهات الأسطر، كما تؤكد
\cref{prop:b2:linalg:transposerank} --- فالعلاقة الوحيدة
``السطر$_3$ = السطر$_1$ + السطر$_2$'' \emph{هي}
مُبيد فضاء الأعمدة. ورتبة الأسطر ($2$ سطرًا حرًّا) ورتبة
الأعمدة تتفقان لا مصادفةً بل لأن كلتيهما تساوي
$\operatorname{rk} A = \operatorname{rk} A^{\mathsf T}$.
\end{example}

\begin{example}[الثنوية تقرأ قاعدة تكميم]\label{ex:b2:linalg:quadratureread}
لماذا توجد قاعدة مثل قاعدة سيمبسون (\cref{exo:b2:linalg:4})
ولماذا هي وحيدة؟ تجيب الثنوية قبل أي حساب.
على الفضاء $E = \R_2[X]$، يكون التكامل $P \mapsto \int_0^1 P$ متجهة
واحدة بعينها في \omterm{def:b2:linalg:dual}{الثنوي} $E^*$ ذي البعد ثلاثة؛ والتقييمات
عند $0$ و$\frac12$ و$1$ تشكّل \emph{أساسًا} للفضاء
$E^*$؛ ومن ثم ينتشر التكامل عليها بكيفية وحيدة --- وهذا
النشر \emph{هو} قاعدة سيمبسون بمعاملاتها. ويعاير عدّ الأبعاد
التوقعات أيضًا: فعلى الفضاء $\R_3[X]$،
لا يمكن عمومًا لثلاثة تقييمات أن تولّد أربعة أبعاد من الصيغ،
ومن ثم فالضبط على كثيرات الحدود من الدرجة الثالثة ليس دينًا
تؤديه الثنوية؛ وأما كون سيمبسون يكامل كثيرات الحدود التكعيبية
بالضبط رغم ذلك فمكافأة تماثل (تلاشي الدرجات الفردية حول
$\frac12$)، ينبغي التحقق منها يدويًا. فالقواعد ذات $n + 1$ عقدة نشرٌ
لصيغة التكامل في أساس من التقييمات للفضاء $\R_n[X]^*$:
ويكلّف الوجود والوحدانية مبرهنة \omterm{def:b2:linalg:dual}{أساس ثنوي} واحدة؛ ولا تكلّف
جهدًا إلا الدرجات الإضافية.
\end{example}

\section{الصيغ المتعددة الخطية المتناوبة}

\begin{definition}\label{def:b2:linalg:alternating}
يكون التطبيق $f \colon E^n \to K$ \emph{$n$-خطي} إذا كان خطيًا في
كل متغيّر، و\emph{متناوبًا}\index{صيغة متناوبة} إذا
انعدم كلما تساوى وسيطان. ويستلزم التناوبُ
\emph{ضدّ التماثل}: فتبديل وسيطين يغيّر الإشارة
(بنشر $f(\dots, x + y, \dots, x + y, \dots) = 0$)؛ وبعبارة أعمّ، من أجل
$\sigma \in \mathfrak{S}_n$،
\[
f(x_{\sigma(1)}, \dots, x_{\sigma(n)}) =
\varepsilon(\sigma)\, f(x_1, \dots, x_n),
\]
وذلك بتفكيك $\sigma$ إلى مبادلات
(\cref{thm:b2:structures:signature}).
\end{definition}

\begin{theorem}[المبرهنة الأساسية للمحددات]\label{thm:b2:linalg:detspace}
ليكن $\dim E = n$ وليكن $\mathcal{B} = (e_1, \dots, e_n)$ أساسًا. بعد
فضاء الصيغ $n$-الخطية المتناوبة على $E$ هو $1$:
فكل صيغة كهذه مضاعف للصيغة
\[
\det{}_{\mathcal{B}}(x_1, \dots, x_n)
= \sum_{\sigma \in \mathfrak{S}_n} \varepsilon(\sigma)
\prod_{i=1}^{n} a_{\sigma(i),\,i},
\qquad
x_j = \sum_{i} a_{ij} e_i ,
\]
و$\det_{\mathcal{B}}$ هي الوحيدة التي تأخذ القيمة $1$ عند
$\mathcal{B}$.\index{محدد!إنشاء المحدد}
\end{theorem}

\begin{proof}
لتكن $f$ صيغة $n$-خطية متناوبة. بنشر كل وسيط على
$\mathcal{B}$ بتعدد الخطية،
\[
f(x_1, \dots, x_n)
= \sum_{i_1, \dots, i_n} a_{i_1,1}\cdots a_{i_n,n}\,
f(e_{i_1}, \dots, e_{i_n}).
\]
تنعدم الحدود ذات الدليل المتكرر (بالتناوب)؛ والمجموعات
$(i_1, \dots, i_n)$ الباقية هي المتباينة، أي $i_k =
\sigma(k)$ من أجل تبديلة $\sigma$، ويُعيد ضدّ التماثل ترتيب
$f(e_{\sigma(1)}, \dots, e_{\sigma(n)}) = \varepsilon(\sigma)
f(e_1, \dots, e_n)$. ومنه
\[
f = f(e_1, \dots, e_n) \cdot \det{}_{\mathcal{B}} :
\]
فكل \omterm{def:b2:linalg:alternating}{صيغة متناوبة} هي ذلك المضاعف، بشرط أن تكون
$\det_{\mathcal{B}}$ نفسها (أي المجموع المعروض) \emph{هي} الصيغة
$n$-الخطية المتناوبة تأخذ القيمة $1$ عند $\mathcal B$.
وتعدد الخطية واضح (فكل حدّ خطي في كل عمود).
والقيمة عند $\mathcal B$: الحدّ غير المعدوم الوحيد هو $\sigma =
\mathrm{id}$. والتناوب: لنفترض أن $x_j = x_k$ (مع $j \neq k$)، بحيث
تحقق أعمدة الإحداثيات $a_{i j} = a_{i k}$ من أجل كل
$i$. نقابل بين كل $\sigma$ و$\sigma' = \sigma\circ(j\,k)$ ---
وهو تقابل ذاتي من الرتبة اثنين بلا نقاط ثابتة على $\mathfrak{S}_n$.
والجداءان المتقابلان متطابقان:
\[
\prod_i a_{\sigma'(i),\,i}
= a_{\sigma(k),\,j}\; a_{\sigma(j),\,k}
\prod_{i \neq j,k} a_{\sigma(i),\,i}
= a_{\sigma(k),\,k}\; a_{\sigma(j),\,j}
\prod_{i \neq j,k} a_{\sigma(i),\,i}
= \prod_i a_{\sigma(i),\,i},
\]
باستعمال تساوي العمودين $j$ و$k$؛ في حين أن
$\varepsilon(\sigma') = -\varepsilon(\sigma)$. فيسهم كل زوج
بالصفر: ومنه ينعدم المجموع.
\end{proof}

\begin{example}[قاعدة ساروس، مستنبطةً ومهدومة]\label{ex:b2:linalg:sarrus}
من أجل $n = 3$ يكون في صيغة التبديلات $3! = 6$ حدود بالضبط.
وبسرد $\mathfrak{S}_3$ حسب الإشارة ---
$\mathrm{id}$ و$(1\,2\,3)$ و$(1\,3\,2)$ زوجية؛
و$(1\,2)$ و$(1\,3)$ و$(2\,3)$ فردية --- نجد
\[
\det A = a_{11}a_{22}a_{33} + a_{21}a_{32}a_{13} +
a_{31}a_{12}a_{23}
- a_{21}a_{12}a_{33} - a_{31}a_{22}a_{13} -
a_{11}a_{32}a_{23} :
\]
وهي بالضبط قاعدة ``الأقطار'' عند ساروس التي تُدرَّس في المدرسة ---
وقد صارت الآن مبرهنة، بإشارات غامضة تبيّن أنها إشارات التبديلات.
وأما الهدم: فمن أجل $n = 4$ هناك $24$
تبديلة، لا يلتقط منها أي مخطط أقطار سوى $8$؛ فليس لقاعدة
ساروس نسخة من الدرجة $4$، ويحلّ محلها النشر بالعوامل
المرافقة (\cref{thm:b2:linalg:detrules}~(4)).
وعدّ الحدود تحذير أيضًا: ففي صيغة التبديلات
$n!$ حدًّا، ومن ثم فهي \emph{تعريف} لا خوارزمية --- إذ
يحسب اختزال الأسطر \omterm{def:b2:linalg:det}{المحدد} $\det$ في $O(n^3)$
عملية بدلًا من ذلك.
\end{example}

\begin{definition}[المحددات]\label{def:b2:linalg:det}
\emph{محدد عائلة}\index{محدد} في أساس هو
$\det_{\mathcal{B}}(x_1, \dots, x_n)$؛ و\emph{محدد المصفوفة}
$A$ هو محدد أعمدتها في الأساس القانوني
--- أي صيغة التبديلات أعلاه؛ و\emph{محدد التشاكل الذاتي}
$u$ هو العدد $\det u$ الذي يحقق
\[
\det{}_{\mathcal{B}}\bigl(u(x_1), \dots, u(x_n)\bigr)
= \det u \cdot \det{}_{\mathcal{B}}(x_1, \dots, x_n)
\quad \text{من أجل كل } x_i
\]
(فالطرف الأيسر صيغة $n$-خطية متناوبة، ومن ثم فهو مضاعف
للصيغة $\det_\mathcal{B}$ حسب \cref{thm:b2:linalg:detspace}؛ والعامل لا
يتعلق بالأساس $\mathcal{B}$).
\end{definition}

\begin{theorem}[حساب المحددات، مبرهَنًا]\label{thm:b2:linalg:detrules}
\begin{enumerate}
  \item $\det(uv) = \det u\,\det v$؛ و$\;\det(AB) = \det A \det B$.
  \item تكون $u$ قابلة للقلب $\iff \det u \neq 0$؛ وتكون عائلة أساسًا
        $\iff$ كان محددها في أساس ما غير معدوم.
  \item $\det(A^{\mathsf T}) = \det A$.
  \item يصحّ النشر بالعوامل المرافقة على أي سطر أو عمود، كما
        ورد في مجلد السنة الأولى؛ وللمصفوفتين المتشابهتين
        \omterm{def:b2:linalg:det}{المحدد} نفسه.
\end{enumerate}
\end{theorem}

\begin{proof}
(1) نطبّق العلاقة المعرِّفة مرتين:
$\det_{\mathcal B}(uv(x_i)) = \det u \cdot \det_{\mathcal
B}(v(x_i)) = \det u \det v \cdot \det_{\mathcal B}(x_i)$.

(2) إذا كان $u$ قابلًا للقلب، فإن $\det u \det u^{-1} = \det \mathrm{id} =
1 \neq 0$. وإلا فالصور $u(e_i)$ مرتبطة؛ وبالتعبير عن إحداها
بدلالة الأخريات ثم النشر، نجد $\det_{\mathcal B}(u(e_i)) = 0$
(فالتناوب يقتل الاتجاهات المتكررة)، ومنه $\det u = 0$. ومحك الأساس
هو القول نفسه من أجل العائلات.

(3) في صيغة التبديلات، نعيد ترقيم كل جداء بالدليل $j =
\sigma(i)$، أي $i = \tau(j)$ مع $\tau = \sigma^{-1}$: فالعوامل
هي الأعداد نفسها بترتيب آخر، ومنه
\[
\prod_{i=1}^{n} a_{\sigma(i),\,i} = \prod_{j=1}^{n}
a_{j,\,\tau(j)} ,
\]
و$\varepsilon(\tau) = \varepsilon(\sigma)^{-1} =
\varepsilon(\sigma)$ (فالقيم $\pm1$؛ و$\varepsilon$
تشاكل). والمجموع على $\sigma$ هو نفسه المجموع على
$\tau$ (فالقلب تقابل على $\mathfrak{S}_n$):
\[
\det A = \sum_{\tau}\varepsilon(\tau)\prod_j a_{j,\tau(j)}
= \det(A^{\mathsf T}),
\]
والمجموع الأخير هو صيغة التبديلات مطبَّقة على
العناصر المنقولة $(A^{\mathsf T})_{ij} = a_{ji}$.

(4) نثبّت العمود $j$ ونُفكك $x_j = \sum_i a_{ij} e_i$
بالخطية: $\det A = \sum_i a_{ij}\, \det(\dots, e_i, \dots)$، ونقلُ
$e_i$ إلى الموضع الأخير ($n - i$ \omterm{def:b2:structures:sn}{مبادلة} للأسطر
و$n - j$ للأعمدة، عبر (3)) يماهي $\det(\dots, e_i, \dots) =
(-1)^{i+j}\Delta_{ij}$ مع الأقصر: وهي بالضبط قاعدة العوامل المرافقة
في السنة الأولى. وأما التشابه: فلدينا $\det(P^{-1}AP) = \det P^{-1}\det A \det P =
\det A$ حسب (1).
\end{proof}

\begin{example}[النشر بالعوامل المرافقة، منفَّذًا]\label{ex:b2:linalg:cofactorwork}
لنحسب
\[
\det\begin{pmatrix}
2 & 1 & 3\\
0 & 4 & 1\\
1 & 2 & 0
\end{pmatrix}
\]
على العمود الأول (بكسلٍ يساوي صفرين: أي واحدًا).
وتتبع الإشارات رقعة الشطرنج $(-1)^{i+j}$:
\[
2\,\det\begin{pmatrix}4 & 1\\ 2 & 0\end{pmatrix}
- 0
+ 1\cdot\det\begin{pmatrix}1 & 3\\ 4 & 1\end{pmatrix}
= 2(0 - 2) + (1 - 12) = -15 .
\]
وللتحقق المتقاطع بقاعدة ساروس (\cref{ex:b2:linalg:sarrus}): $0 + 1 + 0
- 12 - 0 - 4 = -15$. إنها استراتيجية لا عقيدة: انشر على
الخط الأكثر أصفارًا، وإذا لم يكن في أي خط صفر فاصنع
أصفارًا أولًا بعمليات على الأسطر --- فجولة اختزال واحدة تكلّف
أقلّ من طبقتَي عوامل مرافقة.
\end{example}

\begin{example}[محدد بالقواعد]\label{ex:b2:linalg:onesmatrix}
لتكن $J \in \mathcal{M}_n(K)$ المصفوفة التي كل عناصرها واحد ولتكن $a \in K$؛
ونحسب $\det(aI_n + J)$ بالأدوات التي بُرهن عليها للتوّ. فمجموع كل
عمود من أعمدة $aI_n + J$ يُحسب بالطريقة نفسها: نضيف كل الأسطر إلى
الأول (فلا يتغير \omterm{def:b2:linalg:det}{المحدد} --- إذ إن إضافة مضاعف لسطر إلى
سطر آخر تضيف حدًّا باتجاه متكرر، يقتله التناوب).
فيصير السطر الأول $(a + n, a + n, \dots, a + n)$؛ ثم نُخرج العامل
$a + n$ بالخطية في ذلك السطر، ثم نطرح العمود الأول
من كل عمود آخر: فيبقى ما هو مثلثي بقطر
$(1, a, \dots, a)$. ومنه
\[
\det(aI_n + J) = (a + n)\,a^{\,n-1}.
\]
والفكرة الختامية: الجذران $a = 0$ (بتضاعف $n - 1$) و
$a = -n$ يقولان إن للمصفوفة $J$ القيمة الذاتية $0$ بتضاعف $n - 1$
والقيمة الذاتية $n$ مرة واحدة --- أي طيف المصفوفة ذات الرتبة واحد
$J$، قبل فصل كامل (وستجعل \cref{ch:b2:reduction} هذا
منهجيًا).
\end{example}

\begin{example}[محدد بصيغة التبديلات]\label{ex:b2:linalg:permexample}
من أجل مصفوفة كثيرة الأصفار تكون الصيغة عملية بذاتها:
ففي
\[
A = \begin{pmatrix}
0 & a & 0 & 0\\
0 & 0 & b & 0\\
0 & 0 & 0 & c\\
d & 0 & 0 & 0
\end{pmatrix},
\]
التبديلة الوحيدة التي تلتقط عناصر غير معدومة هي \omterm{def:b2:structures:sn}{الدورة} ذات الطول $4$،
أي $\sigma = (1\,2\,3\,4)$، التي ترسل العمود $1 \to$ إلى السطر $4$، وهكذا؛
و$\varepsilon(\sigma) = (-1)^3 = -1$، ومنه $\det A = -abcd$. (وللتحقق
بثلاث مبادلات أعمدة نبلغ مصفوفة قطرية.)
\end{example}

\begin{example}[محدد فاندرموند بصيغة الجداء]\label{ex:b2:linalg:vandermondework}
من أجل العقد $0, 1, 2$ (المستعملة في قواعد التكميم مثل
قاعدة \cref{exo:b2:linalg:4})، يُحسب \omterm{def:b2:linalg:det}{محدد} فاندرموند للفضاء
\cref{exo:b2:linalg:11} بنظرة واحدة:
\[
\det\begin{pmatrix}
1 & 1 & 1\\
0 & 1 & 2\\
0 & 1 & 4
\end{pmatrix}
= (1 - 0)(2 - 0)(2 - 1) = 2 ,
\]
وبالنشر المباشر على العمود الأول: $1\cdot(4 - 2)
= 2$: وهو التوافق نفسه. وعدم الانعدام من أجل عقد مختلفة هو
نظرية الاستكمال كلها في \omterm{def:b2:linalg:det}{محدد} واحد: فصيغ التقييم
$P \mapsto P(a_i)$ أساسٌ للثنوي بالضبط حين يكون هذا \omterm{def:b2:linalg:det}{المحدد}
غير معدوم، أي دائمًا من أجل $a_i$ مختلفة
--- أي \cref{ex:b2:linalg:dualexamples} مقدَّرةً كمًّا.
\end{example}

\section{الأثر، من جديد}

\begin{proposition}\label{prop:b2:linalg:trace}
الأثر $\operatorname{tr} \colon \mathcal{M}_n(K) \to K$ هو الصيغة الخطية
الوحيدة التي تحقق $\operatorname{tr}(AB) =
\operatorname{tr}(BA)$ و$\operatorname{tr}(I_n) = n$ (من أجل
$\operatorname{char} K = 0$)؛ وأثر التشاكل الذاتي معرَّف جيدًا
عبر أي تمثيل مصفوفي، ويكون
\[
\operatorname{tr}(u) = \sum_{i} e_i^*\bigl(u(e_i)\bigr)
\]
في أي أساس --- فالثنوية تكتب الأثر بلا أساس.
\end{proposition}

\begin{proof}
بُرهن على $\operatorname{tr}(AB) = \operatorname{tr}(BA)$ وعلى عدم التعلق
بالأساس في السنة الأولى. أما الوحدانية: فالصيغة الخطية $t$ التي $t(AB) =
t(BA)$ تُعدم كل مبدّل $AB - BA$. وندّعي أن المبدّلات
تولّد الفضاء الفوقي المعدوم الأثر، وبعده $n^2 - 1$. وتكفي
عائلتان من المبدّلات. فقاعدة ضرب المصفوفات الأولية هي
$E_{ab}E_{cd} = \delta_{bc}E_{ad}$. ومن أجل
$i \neq j$ تعطي
\[
E_{ii}E_{ij} - E_{ij}E_{ii} = E_{ij} - 0 = E_{ij}
\]
(فالجداء الثاني $E_{ij}E_{ii} = \delta_{ji}E_{ii} = 0$
لأن $j \neq i$): فكل مصفوفة $E_{ij}$ خارج القطر مبدّل.
و
\[
E_{ij}E_{ji} - E_{ji}E_{ij} = E_{ii} - E_{jj} .
\]
والمصفوفات $E_{ij}$ (مع $i \neq j$، وعددها $n^2 - n$) مع المصفوفات
$E_{11} - E_{jj}$ (مع $j \geq 2$، وعددها $n - 1$) تشكّل $n^2 - 1$
مصفوفة معدومة الأثر ومستقلة خطيًا: فهي تولّد
الفضاء الفوقي $\ker\operatorname{tr}$. ومنه فإن $t$ تنعدم حيث تنعدم
$\operatorname{tr}$ وتتحلل عبرها: أي $t =
c\operatorname{tr}$؛ ثم يفرض $t(I) = n$ أن $c = 1$. وأما الصيغة المعروضة:
فالعنصر القطري رقم $i$ في مصفوفة $u$ هو بالضبط
$e_i^*(u(e_i))$.
\end{proof}

\begin{remark}[مزالق شائعة]\label{rem:b2:linalg:pitfalls}
(1) \omterm{def:b2:linalg:det}{المحدد} $n$-خطي في \emph{الأعمدة}، لا خطي في
المصفوفة: فلدينا $\det(A + B) \neq \det A + \det B$ عمومًا،
و$\det(\lambda A) = \lambda^n\det A$ لا $\lambda\det A$. (2) النقل يعكس الجداءات:
$(vu)^{\mathsf T} = u^{\mathsf T}v^{\mathsf T}$؛ ونسيان
هذا العكس يُفسد كل حساب فيه مقلوبات. (3) \omterm{def:b2:linalg:annihilator}{المُبيد}
$F^\circ$ يقع في $E^*$ لا في $E$: ولا يصير ``المتممة
المتعامدة'' المألوفة إلا بعد أن يماهي جداءٌ سلّمي
$E$ مع $E^*$ (\cref{ch:b2:quadratic})؛ ولا شيء في
هذه المماهاة قانوني. (4) قولُ ``رتبة الأسطر تساوي رتبة
الأعمدة'' لا يعني أن \emph{فضاء} الأسطر يساوي فضاء الأعمدة ---
فهما يقعان في فضاءين مختلفين ($K^n$ و$K^m$) ويرتبطان
عبر \cref{prop:b2:linalg:transposerank} لا بالتساوي.
(5) صيغة التبديلات أداة برهان: أما من أجل الأعداد فاستعمل
عمليات الأسطر والعوامل المرافقة (\cref{ex:b2:linalg:sarrus}).
\end{remark}

\begin{example}[جداء الأثر يفكك فضاء المصفوفات]\label{ex:b2:linalg:tracepairing}
على $\mathcal{M}_2(\R)$ بالجداء $\langle A, B\rangle =
\operatorname{tr}(AB)$ الوارد في \cref{exo:b2:linalg:9}: نفكك $M =
\begin{psmallmatrix}1 & 4\\ 2 & 3\end{psmallmatrix}$ إلى
جزء متماثل وآخر ضدّ متماثل،
\[
M = S + A, \qquad
S = \tfrac12(M + M^{\mathsf T}) =
\begin{pmatrix}1 & 3\\ 3 & 3\end{pmatrix},
\qquad
A = \tfrac12(M - M^{\mathsf T}) =
\begin{pmatrix}0 & 1\\ -1 & 0\end{pmatrix}.
\]
عندئذٍ $\operatorname{tr}(SA) = \operatorname{tr}
\begin{psmallmatrix}-3 & 1\\ -3 & 3\end{psmallmatrix} = 0$: فالجزآن
``متعامدان'' من أجل جداء الأثر --- وهي حالة خاصة من الواقعة
العامة (المبرهَن عليها في مسألة نهاية الأسبوع في هذا الفصل)
القائلة إن المصفوفات ضدّ المتماثلة تشكّل بالضبط مُبيد
المصفوفات المتماثلة. فالثنوية ترى التفكيك
$\mathcal{M}_n = \mathcal{S}_n \oplus
\mathcal{A}_n$ قبل اختيار أي جداء سلّمي.
\end{example}

\begin{remark}[آفاق داخل هذا المجلد]\label{rem:b2:linalg:perspectives}
راقب إنشاءات هذا الفصل الثلاثة وهي تغيّر أزياءها فيما
يأتي. يعود \emph{\omterm{def:b2:linalg:transpose}{المنقول}} في
\cref{ch:b2:reduction}: فالتطبيقان $u$ و$u^{\mathsf T}$ يشتركان في
القيم الذاتية بتضاعفات هندسية متساوية (مسألة نهاية الأسبوع
في هذا الفصل، السؤال 15)، ولهذا لا يختلف تحليل الأسطر عن
تحليل الأعمدة أبدًا. ويصير \emph{\omterm{def:b2:linalg:det}{المحدد}} دالة وسيط في
\cref{ch:b2:reduction} ($\chi_u(X) = \det(X\,\mathrm{id} - u)$) وياكوبيًا في
\cref{ch:b2:multint}، حيث يتحول تعدد خطيته إلى عامل
تغيير المتغيرات. ويبذر \emph{الأثر} لا متغيّرات
التشابه: فهو المعامل الثاني في
$\chi_u$، ومجموع القيم الذاتية، ثم تكاملَ
القطر في متطابقات من نمط \cref{ch:b2:fourier}. فصلٌ واحد
في \omterm{def:b2:structures:algebra}{الجبر} الخطي، وثلاثة ظلال طويلة.
\end{remark}

\begin{remark}[أين يُستعمَل هذا الفصل]
ليس \omterm{def:b2:linalg:dual}{الفضاء الثنوي} تجريدًا لذاته: فالمبيدات والمنقولات تُدير
نظرية قابلية حلّ الجمل الخطية (وتبرهن مسألة نهاية الأسبوع في
هذا الفصل على بديل فريدهولم في البعد المنتهي انطلاقًا منها)،
وتعود الجداءات غير المنحلة على صورة الصيغة القطبية في
\cref{ch:b2:quadratic} والمرافق في
\cref{ch:b2:hermitian}، ويُشغّل \omterm{def:b2:linalg:det}{المحدد} المبني هنا
\cref{ch:b2:reduction} كله. وفي مجلد السنة الثالثة تصير الثنوية
نفسها، منقولةً إلى البعد غير المنتهي، مبرهنةَ التمثيل عند
ريس ونظريةَ فريدهولم على فضاءات هيلبرت --- مع التراص بديلًا
عن عدّ الأبعاد المستعمل هنا.
\end{remark}

\section{تمارين}

\begin{exercise}[$\star$]\label{exo:b2:linalg:1}
في $\R^3$، لتكن $\varphi_1(x,y,z) = x + y$، $\varphi_2 = y + z$،
$\varphi_3 = x + z$. برهن على أن $(\varphi_1, \varphi_2, \varphi_3)$
أساس للفضاء $(\R^3)^*$ وجد الأساس للفضاء $\R^3$ الذي هو ثنويّه.
\end{exercise}

\begin{exercise}[$\star$]\label{exo:b2:linalg:2}
احسب بصيغة التبديلات محددَي المصفوفتين
\[
\begin{pmatrix} 0 & 0 & a\\ 0 & b & 0\\ c & 0 & 0 \end{pmatrix},
\qquad
\begin{pmatrix}
a & b & 0 & 0\\
c & d & 0 & 0\\
0 & 0 & e & f\\
0 & 0 & g & h
\end{pmatrix},
\]
وصُغ قاعدة الكتل القطرية التي توحي بها الثانية.
\end{exercise}

\begin{exercise}[$\star$]\label{exo:b2:linalg:3}
ليكن $F = \{(x,y,z,t) \in \R^4 : x + y = z + t \text{ و} x = 2y\}$.
أعطِ أساسًا للفضاء $F^\circ$ وتحقق من
\cref{thm:b2:linalg:annihilator} على الأبعاد.
\end{exercise}

\begin{exercise}[$\star\star$]\label{exo:b2:linalg:4}
لتكن $a_0, \dots, a_n$ نقاطًا مختلفة من $K$ ولتكن $\varphi_i \colon
P \mapsto P(a_i)$ على $K_n[X]$. برهن على أن $(\varphi_0, \dots,
\varphi_n)$ أساس للفضاء $K_n[X]^*$، وعيّن أساسه القبلي،
وانشر الصيغة $P \mapsto \int_0^1 P(t)\,\dd t$ (من أجل $K = \R$ و
$n = 2$ و$a_i = 0, \frac12, 1$) في هذا الأساس --- ستتعرّف على
قاعدة سيمبسون.
\end{exercise}

\begin{exercise}[$\star\star$]\label{exo:b2:linalg:5}
ليكن $u \in \mathcal{L}(E)$ مع $\dim E = n$ و$\operatorname{rk} u
= 1$. برهن على أن $u = \varphi(\cdot)\, a$ من أجل متجهة $a$ وصيغة
$\varphi$؛ وعلى أن $\operatorname{tr} u = \varphi(a)$؛ وعلى أن
$u^2 = (\operatorname{tr} u)\, u$. استنتج $\det(I + u) = 1 +
\operatorname{tr} u$.
\end{exercise}

\begin{exercise}[$\star\star$]\label{exo:b2:linalg:6}
برهن على أن كل فضاء فوقي من $\mathcal{M}_n(K)$ (مع $n \geq 2$)
يحتوي مصفوفة قابلة للقلب. \emph{إرشاد: الفضاء الفوقي هو $\{M :
\operatorname{tr}(AM) = 0\}$ من أجل $A \neq 0$ ما
(\cref{exo:b2:linalg:9}). فإذا كانت $A$ سلّمية، فأبرِز مصفوفة
قابلة للقلب معدومة الأثر؛ وإلا فجد مصفوفة $M$ قابلة للقلب تجعل
$AM$ ذات قطر معدوم --- وتفي بذلك مصفوفة شبيهة بمصفوفة تبديلة.}
\end{exercise}

\begin{exercise}[$\star\star$]\label{exo:b2:linalg:7}
(مشتق \omterm{def:b2:linalg:det}{المحدد}) من أجل $A \in \mathcal{M}_n(\R)$، برهن انطلاقًا
من تعدد الخطية على أن
\[
\frac{\dd}{\dd t}\Big|_{t=0} \det(I_n + tA) = \operatorname{tr} A ,
\]
واستنتج $\det(\eu^{tA}) = \eu^{t\operatorname{tr} A}$ بافتراض
قابلية $t \mapsto \det(\eu^{tA})$ للاشتقاق وخاصية الزمرة
$\eu^{(s+t)A} = \eu^{sA}\eu^{tA}$ (المثبتة في
\cref{ch:b2:diffeq}).
\end{exercise}

\begin{exercise}[$\star\star$]\label{exo:b2:linalg:8}
(مصفوفة دائرية، $3 \times 3$) ليكن $j = \eu^{2\iu\pi/3}$ ولتكن
\[
C = \begin{pmatrix}
a & b & c\\
c & a & b\\
b & c & a
\end{pmatrix} \in \mathcal{M}_3(\C).
\]
تحقق من أن أعمدة مصفوفة فاندرموند للعناصر $1, j, j^2$ متجهات ذاتية
للمصفوفة $C$، واستنتج
\[
\det C = (a + b + c)(a + bj + cj^2)(a + bj^2 + cj).
\]
\end{exercise}

\begin{exercise}[$\star\star\star$]\label{exo:b2:linalg:9}
برهن على أن كل صيغة خطية $t$ على $\mathcal{M}_n(K)$ تُكتب $M \mapsto
\operatorname{tr}(AM)$ من أجل $A$ وحيدة: أي إن التطبيق $A \mapsto
\operatorname{tr}(A\,\cdot)$ تماثل من
$\mathcal{M}_n(K)$ على ثنويّه. استنتج من جديد قول الوحدانية في
\cref{prop:b2:linalg:trace}.
\end{exercise}

\begin{exercise}[$\star\star\star$]\label{exo:b2:linalg:10}
ليكن $u, v \in \mathcal{L}(E)$ مع $u \circ v - v \circ u = u$.
برهن على أن $u$ عديم القوى. \emph{إرشاد: بيّن أن
$\operatorname{tr}(u^k) = 0$ من أجل كل $k \geq 1$ (احسب $u^k v -
v u^k$ بالتراجع)، ثم استعمل الواقعة التالية، التي يُبرهن عليها
بمتطابقات نيوتن أو بالتراجع على البعد: كل تشاكل ذاتي
لفضاء متجهي على $\C$ تنعدم آثار كل قواه يكون عديم القوى.
واعمل على $\C$.}
\end{exercise}

\begin{exercise}[$\star\star$]\label{exo:b2:linalg:11}
(فاندرموند) من أجل $a_0, \dots, a_n \in K$، برهن على أن
\[
\det\begin{pmatrix}
1 & 1 & \cdots & 1\\
a_0 & a_1 & \cdots & a_n\\
\vdots & \vdots & & \vdots\\
a_0^n & a_1^n & \cdots & a_n^n
\end{pmatrix}
= \prod_{0 \leq i < j \leq n} (a_j - a_i).
\]
\emph{(انظر إلى \omterm{def:b2:linalg:det}{المحدد} بوصفه كثير حدود في $a_n$: عيّن
درجته وجذوره ومعامله الرئيسي؛ ثم استعمل التراجع.)}
\end{exercise}

\begin{exercise}[$\star\star\star$]\label{exo:b2:linalg:12}
ليكن $A, B, C, D \in \mathcal{M}_n(K)$ مع $K$ غير منته، وليكن
$CD = DC$. برهن على أن
\[
\det\begin{pmatrix} A & B\\ C & D\end{pmatrix}
= \det(AD - BC).
\]
\emph{(عالج أولًا حالة $D$ قابلة للقلب، بالضرب من اليمين في
$\begin{psmallmatrix} I & 0\\ -D^{-1}C & I\end{psmallmatrix}$؛ ثم عوّض $D$ بالمصفوفة $D + tI$ وقارن بين كثيري حدود في $t$.)}
\end{exercise}

\section{مسألة: بديل فريدهولم}

متى تقبل الجملة الخطية $u(x) = b$ حلًا؟ الجواب
الكامل قولٌ في الثنوية: \emph{بالضبط حين تُبيد $b$
كلُّ صيغة خطية تُبيد صورة $u$} --- وتلك الصيغ قابلة
للحساب، إذ هي نواة \omterm{def:b2:linalg:transpose}{المنقول}. وتبني مسألة نهاية الأسبوع هذه
معجم الثنوية الكامل في البعد المنتهي (تحليل الصيغ، والثنوية
المزدوجة، وحساب المبيدات، \omterm{def:b2:linalg:transpose}{والمنقول})، وتبرهن على
\emph{بديل فريدهولم} في البعد المنتهي، وتُختم بصيغة الأثر
وبتوصيف: الأثر هو اللامتغيّر الخطي الوحيد للتشابه. وفي كل ما
يلي، $E$ و$F$ فضاءان متجهيان منتهيا البعد على $K$،
مع $n = \dim E$.

\begin{problem}[{مسألة نهاية الأسبوع --- الثنوية في البعد المنتهي
وبديل فريدهولم}]
\label{pb:b2:linalg:1}
ترميز: من أجل $S \subseteq E^*$، يكون \emph{\omterm{def:b2:linalg:annihilator}{المبيد} القبلي} هو
$S_\circ = \{x \in E : \varphi(x) = 0 \text{ من أجل كل } \varphi
\in S\}$؛ والمبيدات $F^\circ$ والمنقولات $u^{\mathsf T}$
هي مبيدات ومنقولات \cref{def:b2:linalg:annihilator} و
\cref{def:b2:linalg:transpose}.

\medskip
\noindent\textbf{الجزء الأول --- مبرهنة التحليل المساعدة.} ليكن
$\varphi_1, \dots, \varphi_p, \varphi \in E^*$.
\begin{enumerate}
  \item ليكن $\Phi \colon E \to K^p$، $x \mapsto (\varphi_1(x),
        \dots, \varphi_p(x))$. عيّن $\ker\Phi$، وبيّن
        أن $\Phi^{\mathsf T}$ يرسل صيغ إحداثيات $K^p$ إلى
        العناصر $\varphi_i$، واستنتج
        \[
        \dim \bigl(\ker\varphi_1 \cap \dots \cap
        \ker\varphi_p\bigr) = n - \dim
        \operatorname{Vect}(\varphi_1, \dots, \varphi_p).
        \]
  \item (مبرهنة التحليل المساعدة) برهن على التكافؤ:
        \[
        \varphi \in \operatorname{Vect}(\varphi_1, \dots,
        \varphi_p)
        \iff
        \ker\varphi_1 \cap \dots \cap \ker\varphi_p \subseteq
        \ker\varphi .
        \]
  \item استنتج: تكون $(\varphi_1, \dots, \varphi_p)$ حرة إذا وفقط إذا
        كان بعد $\bigcap_i \ker\varphi_i$ يساوي $n - p$؛ وأن
        كل فضاء جزئي بعده المرافق $p$ تقاطعُ $p$
        فضاء فوقي، لا أقلّ.
  \item في $\R^4$، لتكن $\varphi_1 = x + y - z$، $\varphi_2 = y +
        z - t$، $\psi = x + 2y - t$ ولتكن $\psi' = x + y + t$.
        قرّر، بمبرهنة التحليل المساعدة، هل تنتمي $\psi$ و
        $\psi'$ إلى $\operatorname{Vect}(\varphi_1,
        \varphi_2)$.
  \item على $E = \R_2[X]$، بيّن أن $\psi_0 \colon P \mapsto
        P(0)$ و$\psi_1 \colon P \mapsto P(1)$ و$\psi_2 \colon P
        \mapsto \int_0^1 P(t)\dd t$ تشكّل أساسًا للفضاء $E^*$،
        واحسب الأساس $(P_0, P_1, P_2)$ للفضاء $E$ الذي هو
        ثنويّه، وجد $P \in \R_2[X]$ الوحيدة التي $P(0)
        = 1$ و$P(1) = 2$ و$\int_0^1 P = \frac32$.
\end{enumerate}

\medskip
\noindent\textbf{الجزء الثاني --- الثنوية المزدوجة وحساب
المبيدات.}
\begin{enumerate}[resume]
  \item بيّن أن \emph{تطبيق التقييم} $J \colon E \to
        E^{**}$، $J(x)(\varphi) = \varphi(x)$، خطي
        ومتباين، ومن ثم فهو تماثل في البعد المنتهي.
  \item (\omterm{def:b2:linalg:annihilator}{المبيد} المزدوج) بيّن أن $J(F) = F^{\circ\circ} :=
        (F^\circ)^\circ$ من أجل كل فضاء جزئي $F \subseteq E$: أي إن
        \omterm{def:b2:linalg:annihilator}{مبيد} \omterm{def:b2:linalg:annihilator}{المبيد}، بالمماهاة $J$، هو الفضاء الجزئي نفسه.
  \item برهن على حساب المبيدات: $(F + G)^\circ = F^\circ
        \cap G^\circ$ و$(F \cap G)^\circ = F^\circ + G^\circ$.
  \item استنتج (ثم أعد البرهان مباشرة): كل صيغتين غير معدومتين
        لهما النواة نفسها متناسبتان.
  \item (الأساس القبلي) بيّن أنه من أجل كل أساس $(\varphi_1,
        \dots, \varphi_n)$ للفضاء $E^*$ يوجد أساس وحيد
        $(u_1, \dots, u_n)$ للفضاء $E$ يحقق $\varphi_i(u_j) =
        \delta_{ij}$.
\end{enumerate}

\medskip
\noindent\textbf{الجزء الثالث --- حساب المنقولات.}
\begin{enumerate}[resume]
  \item بيّن أن $u \mapsto u^{\mathsf T}$ تقابل خطي
        من $\mathcal{L}(E, F)$ على $\mathcal{L}(F^*, E^*)$،
        وأن $(u^{-1})^{\mathsf T} = (u^{\mathsf T})^{-1}$
        حين تكون $u$ قابلة للقلب.
  \item (الطبيعية) بيّن أن $u^{\mathsf T\mathsf T} \circ J_E
        = J_F \circ u$: أي إن \omterm{def:b2:linalg:transpose}{المنقول}
        المزدوج، بتماثلات التقييم، \emph{هو} $u$.
  \item بيّن أن: $u$ شامل إذا وفقط إذا كان $u^{\mathsf T}$ متباينًا؛
        وأن $u$ متباين إذا وفقط إذا كان $u^{\mathsf T}$ شاملًا.
  \item من أجل $u \in \mathcal{L}(E)$: يكون الفضاء الجزئي $F$ مستقرًا
        تحت $u$ إذا وفقط إذا كان $F^\circ$ مستقرًا تحت
        $u^{\mathsf T}$.
  \item بيّن أن $\ker(u^{\mathsf T} - \lambda\,
        \mathrm{id}_{E^*}) = \bigl(\operatorname{im}(u - \lambda\,
        \mathrm{id}_E)\bigr)^\circ$، واستنتج أن $u$ و
        $u^{\mathsf T}$ لهما القيم الذاتية نفسها بالتضاعفات
        الهندسية نفسها.
\end{enumerate}

\medskip
\noindent\textbf{الجزء الرابع --- بديل فريدهولم.}
\begin{enumerate}[resume]
  \item برهن على أن $\operatorname{im} u = (\ker u^{\mathsf
        T})_\circ$ من أجل $u \in \mathcal{L}(E, F)$، واستنتج
        \emph{بديل فريدهولم} في البعد المنتهي: تقبل
        المعادلة $u(x) = b$ حلًا إذا وفقط إذا كانت كل
        $\psi \in F^*$ تحقق $u^{\mathsf T}\psi = 0$ تحقق أيضًا
        $\psi(b) = 0$.
  \item الصيغة المصفوفية: من أجل $A \in \mathcal{M}_{m,n}(K)$ و$b \in
        K^m$، تتحقق حالة واحدة بالضبط مما يلي: (1) تقبل $Ax = b$
        حلًا؛ (2) توجد $y \in K^m$ تحقق
        $A^{\mathsf T}y = 0$ و$y^{\mathsf T}b = 1$. برهن على
        ``واحدة على الأكثر'' وعلى ``واحدة على الأقل''.
  \item جد كل $b \in \R^3$ التي تقبل من أجلها الجملة
        \[
        x + y = b_1, \qquad y + z = b_2, \qquad x + 2y + z = b_3
        \]
        حلًا، وذلك بحساب نواة المصفوفة المنقولة.
  \item (مسألة نويمان المتقطعة) على $E = \R^n$ (مع $n \geq 3$)،
        نعرّف $L$ بالعلاقة $(Lx)_k = x_k - \frac12(x_{k-1} +
        x_{k+1})$، بأدلة بترديد $n$. بيّن أن $L^{\mathsf T} = L$
        (بالمماهاة القانونية)، وبيّن أن $\ker L$ هو مستقيم
        المتجهات الثابتة \emph{(انظر إلى إحداثية أعظمية)}،
        واستنتج: تقبل $Lx = b$ حلًا إذا وفقط إذا كان $\sum_k b_k = 0$.
\end{enumerate}

\medskip
\noindent\textbf{الجزء الخامس --- صيغة الأثر ومبرهنة
اللاتغيّر.} تذكّر من \cref{exo:b2:linalg:9} أن $A \mapsto
\operatorname{tr}(A\,\cdot)$ يماهي $\mathcal{M}_n(K)$ مع
ثنويّه. ولنفترض $\operatorname{char} K = 0$ (مثلًا $K = \Q, \R,
\C$).
\begin{enumerate}[resume]
  \item بهذه المماهاة، بيّن أن \omterm{def:b2:linalg:annihilator}{مبيد} الفضاء الجزئي
        $\mathcal{S}_n$ للمصفوفات المتماثلة هو الفضاء
        الجزئي $\mathcal{A}_n$ للمصفوفات ضدّ المتماثلة، والعكس
        بالعكس.
  \item بيّن أن \omterm{def:b2:linalg:annihilator}{مبيد} الفضاء الفوقي
        $\mathfrak{sl}_n = \{M : \operatorname{tr} M = 0\}$ هو
        المستقيم $K I_n$؛ أي بعبارة مكافئة، أن كل صيغة خطية تنعدم
        على المصفوفات المعدومة الأثر هي مضاعف للأثر.
  \item بيّن أن كل مصفوفة من $\mathcal{M}_n(K)$ مجموع
        مصفوفتين قابلتين للقلب.
  \item (الأثر هو اللامتغيّر الخطي الوحيد للتشابه) لتكن
        $t$ صيغة خطية على $\mathcal{M}_n(K)$ تحقق
        $t(PMP^{-1}) = t(M)$ من أجل كل $M$ وكل $P$
        قابلة للقلب. بيّن أولًا أن $t(PX) = t(XP)$ من أجل $P$ قابلة للقلب، ثم
        أن $t(BX) = t(XB)$ من أجل \emph{كل} $B$، واستنتج أن $t = c
        \operatorname{tr}$ من أجل $c \in K$ ما.
  \item بيّن أن $\operatorname{rk} u \leq r$ إذا وفقط إذا كان $u$
        مجموع $r$ تطبيقًا رتبة كل منها $\leq 1$، أي $u =
        \sum_{i=1}^{r} \psi_i(\cdot)\,f_i$ مع $\psi_i \in E^*$
        و$f_i \in F$؛ واستنتج $\operatorname{rk}(u + v) \leq
        \operatorname{rk} u + \operatorname{rk} v$.
  \item (تركيب) اجمع المعجم المبرهَن عليه في هذه المسألة:
        الفضاءات الجزئية في مقابل المبيدات، والمجاميع في
        مقابل التقاطعات، والتطبيقات في مقابل المنقولات، وقابلية
        الحلّ في مقابل التعامد مع نواة \omterm{def:b2:linalg:transpose}{المنقول}، والأثر في مقابل
        التشابه. واذكر من أجل كل مدخل السؤال الذي برهن عليه،
        وقل بجملة واحدة ما الذي يحلّ محلّ عدّ الأبعاد حين يصير
        البعد غير منته (ويوضّح مجلد السنة الثالثة ذلك
        على فضاءات هيلبرت).
\end{enumerate}
\end{problem}
